\subsection{Limitations}\label{sec:limitations}

\textsc{TensorGuard} v5 inherits the scoped character of v4 and
extends it; the scope is now explicit in the refinement type
system rather than spread across two passes.

\paragraph{Analyzed Python subset.} We analyze the typed AST
fragment $\mathcal F$ of \verb|nn.Module| classes:
\verb|__init__| with statically resolvable layer constructors,
\verb|forward| with straight-line composition, refinement-typed
helpers, and the operator catalogue of
\Cref{app:ops-extended}. Decorator metaprogramming, dynamic
attribute assignment after \verb|__init__|, and class factories
that synthesize \verb|forward| at runtime are out of scope. Of
the cited benchmark blocks, the in-fragment rate is reported
per-source in \Cref{sec:eval}; out-of-fragment blocks receive an
\textbf{Abstain} verdict, never a false claim.

\paragraph{Truly dynamic control flow.} The calculus assumes
shape- and grad-flag-determinism: the trace structure of a
\verb|forward| call is a function of the static refinement of its
inputs, not of their values. Data-dependent branches
(\verb|if x.sum() > 0|), data-dependent loops, and Python-level
exceptions inside \verb|forward| trigger sound abstention on the
affected sub-tree. This is a strict design choice: the
Dynamo-guard correspondence (\Cref{thm:dynamo-corr}) is stated
modulo Dynamo graph breaks, and abstention is exactly what
\textsc{TensorGuard} emits where Dynamo would have broken the
graph.

\paragraph{Backward verifier scope.} The backward analysis
(\Cref{sec:impl-backward}) tracks the grad-flag refinement
through the operator catalogue and detects severed-tape
conditions; it does \emph{not} verify numerical correctness of
gradients, does not reason about second-order or higher
derivatives, and abstains on custom \verb|autograd.Function|
subclasses (their \verb|backward| body is opaque to the static
analysis). \verb|torch.no_grad()| and \verb|torch.enable_grad()|
scopes are tracked syntactically; runtime context-manager
swapping via \verb|setattr| is out of scope.

\paragraph{Lean parity is operator-rule-only.} The Lean~4
development (\Cref{sec:lean-theorems}) machine-checks the
per-operator transfer-function and refinement-implication lemmas
for $20$ operators. It does \emph{not} mechanise the analyzer
implementation, the AST extractor, or the assume/guarantee
composition rule; soundness of those layers reduces to the Lean
operator lemmas via the pen-and-paper proof of
\Cref{sec:soundness-proof}. The boundary between machine-checked
and pen-and-paper soundness is stated explicitly there and is
unchanged from v4 in character but expanded in coverage.

\paragraph{Dynamo correspondence is one-way at the boundary.} The
correspondence theorem (\Cref{thm:dynamo-corr}) is exact only on
the intersection of \textsc{TensorGuard}'s in-fragment subset and
Dynamo's traceable subset. Outside that intersection,
\textsc{TensorGuard} abstains and Dynamo either installs a
non-shape guard or breaks the graph; we do not claim agreement
on those rows and they are reported separately in
\Cref{sec:eval}.
